% Generated by Claude Fable 5.1 (Anthropic), 2026-09-15 — captions for WIP-53 figures

% ---------- Fig. 1 (timeline figure, final) ----------
\caption{Three timelines of one run, on a small example: a multi-threaded
executor EX1 dispatching callbacks c1--c3 onto threads t11 and t13, the OS
scheduling those threads, and the GPU running the work c2 submits.
\texttt{ros2\_tracing} sees each callback's start and end on its thread; the
OS scheduler sees a thread leave its core, preempted or blocked, but not what
it was running nor what it waits for: t11 is preempted under c1 and c3, t13
is blocked from its submissions until its copies and kernels complete at
$t_5$. \texttt{nsys} sees the copies and kernels and the thread that
submitted them, t13~\textcircled{\scriptsize 1}; which callback that thread
was running at the submission instant, c2, only \texttt{ros2\_tracing}
knows~\textcircled{\scriptsize 2}. To the executor, c2 is one execution from
$t_3$ to $t_5$; where its time went is on no single timeline.}

% ---------- Fig. 2 (three panels: a graph, b layers, c GPU) ----------
\caption{The FISH approach. (a)~The fully expanded FISH graph: the
containment hierarchy from the host down to the individual callbacks, colour
marking the namespace a vertex belongs to and the edges being the containment
relation. (b)~How edges move up the layers. The edges marked
\textcircled{\scriptsize 2} carry the very same message, drawn again at every
layer above. Those marked \textcircled{\scriptsize 3} are one output of a
join: either member callback may publish it, whichever completes the set; at
\textcircled{\scriptsize 4}, where both endpoints are nodes, they merge into
one edge. \texttt{join} and \texttt{polled-msg} edges have both endpoints
inside one node: they exist at layers 3 and 4 and are contained by the node
above, which expanding it brings back. (c)~For a GPU-submitting callback FISH
keeps every execution as a typed DAG, groups the executions into shapes, and
folds the shapes into one conditional graph: a skeleton common to all
executions plus the branches, each with the share of executions that take
it.}
